\begin{abstract}
The GENIUS Act, signed into law on 18~July 2025, requires US payment-stablecoin issuers to hold reserves of cash, deposits, overnight repurchase agreements, government money market funds and Treasury securities maturing within 93~days. We ask whether the association between stablecoin flows and US Treasury-bill yields changed after enactment. Using daily supply of \nCoins{} fiat-backed USD stablecoins and secondary-market yields from \sampleStart{} to \sampleEnd{}, we estimate local projections with a post-enactment interaction, following a pre-analysis plan committed before any regression was run. Before enactment, a USD~10\,bn five-day inflow is associated with a \bPre{}\,bp change in the 3-month yield after ten business days (90\% CI \bPreLo{} to \bPreHi{}). After enactment the estimate is \bPost{}\,bp (90\% CI \bPostLo{} to \bPostHi{}); the difference, \bDiff{}\,bp (\bDiffP{}), has the opposite sign to our hypothesis and is not statistically distinguishable from zero. We find no post-enactment shift of the response towards the shortest maturities and no post-enactment dominance of outflows over inflows. A placebo break in July~2024 is larger and significant (\rbPlDiff{}\,bp, \rbPlDiffP{}), so the relationship is unstable over time and no break can be attributed to the Act. The reserve rules bind only from January~2027 and the post-enactment sample is short, so the design has little power: the null is weak evidence of no change. Data, code and every number in this paper are reproducible from \repo.
\end{abstract}

\begin{IEEEkeywords}
stablecoins, Treasury bills, GENIUS Act, local projections, safe-asset demand, pre-registration
\end{IEEEkeywords}
